% Reference entries, C (AN-11).  Law 1: os/iword/app/face_a.c prompt_key,
% format_key, block_key, retrieve_key (Cancel F1 / Esc in every prompt),
% paint_status (POS with Caps Lock), cursor_key (the arrows, End and the
% Home prefix), os/iword/core/edit.c (the unit walk: a code is a cursor
% position of its own; iw_ed_line_home / line_start / doc_top / doc_end /
% to_line), codes.c (the "Delete [code]? No (Yes)" question, the partner
% deleted with an attribute code; codes InitechWord does not create are
% stepped over).  Shots: ../wordperfect51-decomp mint MINT11, MINT27.
% GATES: test-pubs-iword test-iword-codes test-iword-core test-iword-face test-iword-reveal-codes
\Gates{test-pubs-iword test-iword-codes test-iword-core test-iword-face test-iword-reveal-codes}
\begin{Feature}{Cancel}
Cancel backs out of a menu, prompt or question without doing anything.  You
can press it whenever the status line shows a menu or asks you something.
On the editing screen, with nothing to cancel, the same key restores
deleted text (see \InRef{Undelete}).

To cancel,
\begin{WPSteps}
\item Press \FK{Cancel}{F1}
\Or
Press \K{Esc}.
\end{WPSteps}
\Documents{wpkey}{F1}
\Documents{wpkey}{Esc}

\Notes
\Topic{Block}
While Block is on, Cancel turns Block off and leaves the text as it was.

\Topic{Error Messages}
When the status line shows a message that begins \Scr{ERROR:}, any key --
Cancel included -- takes you back to the prompt where the error happened,
so that you can try again.  Press Cancel once more to return to the
document.

\Topic{Exit and Save}
Cancel at any of the questions of Exit (\Scr{Save Document?},
\Scr{Exit WP?}) or at the name prompt of Save or Retrieve returns you to the
document as it was.

\Topic{Format Screens}
On the Format screens Cancel goes back one screen at a time: from a
measurement you are typing to the Format: Line screen (leaving the setting
as it was), from Format: Line to Format, and from Format to the document.
\FK{Exit}{F7} on the Format: Line screen returns straight to the document.

\Topic{Moving Text}
At \Scr{Move cursor; press Enter to retrieve.} Cancel leaves the text where
it is now: text you moved stays out of the document, but you can still put
it back with Move's Retrieve (see \InRef{Move, Retrieve}).  When InitechWord
has highlighted a sentence, paragraph or page for Move, Cancel removes the
highlight.

\Topic{On the Editing Screen}
With no menu or prompt on the screen, \FK{Cancel}{F1} starts Undelete: it
shows the last text you deleted and offers to restore it (see
\InRef{Undelete}).  \K{Esc} does nothing on the editing screen.

\Topic{Search and Replace}
Cancel at the \Scr{Srch:} prompt leaves the cursor where it was.  At the
questions and strings of Replace it leaves Replace; at
\Scr{Replace? Yes (No)} it stops replacing.  See \InRef{Search} and
\InRef{Replace}.

\Topic{Setup}
Cancel (or Esc) on any Setup screen returns to the document; Exit (F7)
goes back one screen.  See \InRef{Setup}.

\SeeAlso{Block; Exit; Format; Move, Retrieve; Undelete}
\end{Feature}

\begin{Feature}{Caps Lock}
The Caps Lock key of your keyboard works in InitechWord as in every FLAIR
application.  While Caps Lock is on, the status line shows the word
\Scr{Pos} in capitals, \Scr{POS}, as a reminder.

To see whether Caps Lock is on,
\begin{WPSteps}
\item Look at the status line: \Scr{POS} means Caps Lock is on, \Scr{Pos}
that it is off.
\end{WPSteps}

\BeginReplay{ref-caps}
\CapsLock{on}\Keys{End}%
With Caps Lock on, a new document shows this status line:
\WordScreen[rows=1]{ref-caps}{The status line with Caps Lock on}
\CapsLock{off}%
\EndReplay

\SeeAlso{Status Line}
\end{Feature}

\begin{Feature}{Codes}
InitechWord keeps everything that formats a document as hidden codes in the
text.  When you press Bold, a hard return, a tab or a margin setting, a code
goes into the document at the cursor.  The editing screen shows what the
codes do, but not the codes; Reveal Codes shows them (see \InRef{Reveal
Codes}).  Appendix B lists every code and its meaning.

Most features put one code in the document.  Appearances and sizes put a
pair of codes, one where the attribute begins and one where it ends: the
\emph{on} code is named in capitals (\Code{[BOLD]}), the \emph{off} code in
lower case (\Code{[bold]}).

To delete a code,
\begin{WPSteps}
\item Press \FK{Reveal Codes}{Alt-F3} so that you can see the code.
\item Move the cursor onto the code (the code is highlighted) and press
\K{Delete}
\Or
Move the cursor just after the code and press \K{Backspace}.
\item Press \FK{Reveal Codes}{Alt-F3} again to return to the editing screen.
\end{WPSteps}
\begin{WPNote}Deleting either code of a pair deletes both.\end{WPNote}

\Notes
\Topic{Codes InitechWord Keeps but Does Not Make}
A document written by WordPerfect 5.1 can hold codes that InitechWord does
not put in documents itself: \Code{[Center]}, \Code{[Flsh Rgt]} and
\Code{[ ]} (a hard space), and codes Reveal Codes shows as \Code{[?]}.
InitechWord keeps them where they are and saves them unchanged.  Backspace
and Delete step over them instead of deleting them, in Reveal Codes as on
the editing screen.  The editing screen draws the text after a
\Code{[Center]} or \Code{[Flsh Rgt]} at the left margin.

\Topic{Cursor on a Code}
Each code is a place the cursor can stop.  On the editing screen a code takes
no room, so pressing \K{Right} across one may seem to leave the cursor where
it was; in Reveal Codes you see it move onto the code.

\Topic{Deleting Codes on the Editing Screen}
If Backspace or Delete reaches a code of an appearance or size, a margin
setting, a justification or a line spacing while Reveal Codes is off,
InitechWord does not delete it at once.  It asks, naming the code as Reveal
Codes shows it:
\StatusLine{Delete [BOLD]? No (Yes)}
Type \Typed{y} to delete the code (and, for a pair, its partner).  Any other
key answers No: nothing is deleted, and the cursor moves past the code -- to
before it for Backspace, after it for Delete -- so that the next Backspace
or Delete reaches the text.  Tabs, indents, hard returns and hard pages are
deleted without a question.

\BeginReplay{ref-codes}
For example, type \TypeRun{Read the }(with the space), press \Press{Bold}{F6}
and type \TypeRun{new}.  The cursor is just before \Code{[bold]}.  Press
\PressKey{Delete}:
\WordScreen[rows=1]{ref-codes-1}{Delete meets the off code of bold}
Type \TypeName{y}: both bold codes are deleted and \textit{new} is ordinary
text again.
\WordScreen[rows=1]{ref-codes-2}{The bold codes deleted}
\EndReplay

\Topic{Order of Codes}
When you choose a format at the cursor, InitechWord puts its code there.
Margin settings and justification take effect at the start of a line, so if
the cursor is not at the left margin InitechWord first ends the line there
(see \InRef{Margins, Left and Right} and \InRef{Justification}).

\SeeAlso{Appearance; Delete; Reveal Codes; Size}
\end{Feature}

\begin{Feature}{Cursor Movement}
The cursor keys move the cursor through the text without changing it.  The
cursor marks where the next character you type goes; it is the cell drawn
in reverse video on the editing screen.

To move the cursor,
\begin{WPSteps}
\item Press an arrow key to move one character left or right, or one line
up or down.
\item Press \K{End} to move to the end of the line.
\item Press \K{Home} before an arrow key to move farther (see the table
below).
\end{WPSteps}
\Documents{wpkey}{Up}
\Documents{wpkey}{Down}
\Documents{wpkey}{Left}
\Documents{wpkey}{Right}
\Documents{wpkey}{End}
\Documents{wpkey}{Home}
\Documents{wpkey}{Home,Left}
\Documents{wpkey}{Home,Right}
\Documents{wpkey}{Home,Up}
\Documents{wpkey}{Home,Down}
\Documents{wpkey}{Home,Home,Up}
\Documents{wpkey}{Home,Home,Down}
\Documents{wpkey}{Home,Home,Home,Left}
\Documents{wpkey}{Home,Home,Home,Up}

\K{Home} is a prefix key: pressing it does nothing by itself, but changes
what the next arrow key does.  Press Home, release it, then press the arrow.
\par\addvspace{4pt}
\noindent\begin{tabular}{@{}lp{50mm}@{}}
\textbf{Keys} & \textbf{Moves the cursor to}\\
\K{Left}, \K{Right} & the next character to the left or right\\
\K{Up}, \K{Down} & the line above or below\\
\K{End} & the end of the line\\
\K{Home, Left} & the start of the line (after any codes there)\\
\K{Home, Home, Left} & the start of the line (after any codes there)\\
\K{Home, Home, Home, Left} & the start of the line, before any codes there\\
\K{Home, Right} & the end of the line\\
\K{Home, Up} & the top line on the screen\\
\K{Home, Down} & the bottom line on the screen\\
\K{Home, Home, Up} & the top of the document, after the formatting codes there\\
\K{Home, Home, Home, Up} & the top of the document, before any codes\\
\K{Home, Home, Down} & the end of the document\\
\end{tabular}
\par\addvspace{4pt}

\Notes
\Topic{Codes}
The cursor stops on each code as it does on each character.  On the editing
screen a code takes no room, so a press of Left or Right that moves onto or
off a code seems to do nothing; with Reveal Codes on you can see where the
cursor is (see \InRef{Codes}).

\Topic{Home Twice}
Two Homes before Left do what one does, and so do two Homes before Right:
there is no farther left or right than the start and end of a line.  With Up and Down, two Homes go to the top or end of
the whole document.

\Topic{Home, Up and Home, Down at the Edge of the Screen}
If the cursor is already on the top line of the screen, \K{Home, Up} moves
it up one screen; on the bottom line, \K{Home, Down} moves it down one
screen.

\Topic{Moving Up and Down}
Up and Down keep the cursor in the same column as it moves from line to
line.  On a shorter line the cursor goes to the end of the line, and returns
to its column on the next line long enough to have one.

\Topic{Shift, Ctrl and Alt}
The cursor keys do the same with Shift, Ctrl or Alt held as without
(Appendix F).

\Topic{The Screen}
The editing screen shows 23 lines.  When the cursor moves past the top or
bottom of the screen, the screen moves with it.

\SeeAlso{Block; Codes; Status Line}
\end{Feature}
